\chapter{고유기하와 텐서의 대수}

텐서 표기법은 생각의 일부를 대신해 줄 때에만 유용하다. 잘 설계된 표기법은 기계 부품 상자처럼 작동한다. 위 첨자는 아래 첨자와 짝을 이루고, 종류가 다른 대상은 서로 더할 수 없으며, 적법한 축약을 하면 남은 첨자만으로 그 대상의 변환 법칙을 알 수 있다. 이 장치의 직접적인 목적은 기하학적이다. 진정한 곡률과 좌표계가 휘어 있기 때문에 나타나는 곡률을 구분하려 한다.

\section{내부에서 본 곡률}

평범한 종이 한 장에서 시작해 보자. 종이는 늘이거나 찢지 않고 원통 모양으로 말 수 있다. 주변 공간의 관찰자에게 종이는 분명히 휘어져 있다. 그런데 종이 위에서 측정한 거리는 변하지 않는다. 외재적 곡률은 달라졌지만 내재적 기하학은 달라지지 않았다.

이제 종이 위에 갇혀 있는 작은 생명체를 상상해 보자. 이 생명체는 주변 공간을 들여다볼 수 없다. 길이를 재고, 이웃한 방향들을 비교하고, 삼각형을 만들고, 측정 막대를 여기저기로 옮길 수는 있다. 기하학에 필요한 것은 이 내부 실험뿐이어야 한다. 중력에서 이에 대응하는 질문은 어떤 장이 진정한 시공간 곡률을 나타내는지, 아니면 평탄한 시공간을 곡선좌표로 기술한 것에 불과한지 하는 것이다.

이러한 내재적 관점을 구체화하는 한 방법은 곡면을 점과 연결선의 망으로 바꾸는 것이다. 이웃한 점들을 연결하고, 각 연결선의 길이를 기록한 다음, 이음매가 회전하도록 둔다. 정삼각형 격자는 평평하게 놓여 있다. 이음매를 구부려도 격자를 길이 하나 바꾸지 않고 언제든 평면으로 되돌릴 수 있으므로, 기록된 기하학은 달라지지 않는다.

이제 주변 연결선은 그대로 둔 채 한 꼭짓점에 만나는 여섯 연결선의 길이를 늘려 보자. 중심 꼭짓점은 정해진 길이를 어기지 않고는 더 이상 평면에 머물 수 없다. 중심 꼭짓점은 위로 솟아올라 볼록한 부분을 만들어야 한다. 이 곡률을 알아내는 데 바깥에서 보는 관점은 필요하지 않았다. 서로 양립할 수 없는 이웃 간 거리에 이미 곡률이 담겨 있었다.

\begin{figure}[H]
\centering
\includegraphics[width=\linewidth]{lecture_02_figure_01.png}
\caption{내재적 기하학을 나타내는 삼각형 연결선 모형. 한 꼭짓점 주변의 연결선을 바꾸면 길이를 바꾸지 않고는 평평하게 만들 수 없는 볼록한 부분이 생긴다.}
\lectureframe{00:06:40}
\end{figure}

\begin{classroomqa}
\audiencequestion{평탄한 공간의 정확한 정의는 무엇인가요?}
\lecturerresponse{조작적 관점에서 답하면, 그 공간의 내재적 거리 관계를 왜곡 없이 유클리드 배치로 구현할 수 있는 공간이 평탄하다. 아래에서 전개할 언어로 말하면, 관심 영역 전체에서 계량이 상수 크로네커 델타가 되는 좌표를 찾을 수 있을 때 리만 공간은 평탄하다. 한 점에서는 언제나 국소적으로 이렇게 할 수 있다. 이를 어떤 영역 전체에서 할 수 있는지가 자명하지 않은 조건이다. 곡률 텐서는 결국 특별한 좌표 선택에 기대지 않고 이 조건을 표현하게 된다.
\lecturetimestamp{00:09:15}}
\end{classroomqa}

이 답은 아직 정당화하지 않은 계량을 사용한다. 그러므로 잠시 곡률에서 텐서 해석으로 방향을 돌린다. 이 우회는 선택 사항이 아니다. 좌표 의존성과 진정한 기하학을 구분하려면 먼저 계량이 텐서임을 확인해야 한다.

\section{장과 좌표 기술}

일반상대성이론에서 텐서는 대개 텐서 \emph{장}이다. 즉, 각 점에 텐서가 하나씩 있고 그 성분은 점에 따라 변한다. 스칼라장은 가장 간단한 예다. 곡면의 온도는 각 점 \(S(P)\)에서 하나의 수 \(P\)로 나타낼 수 있다. 두 사람이 그 점에 서로 다른 좌표를 붙일 수는 있지만, 그 점의 온도에는 동의해야 한다.

\begin{figure}[H]
\centering
\includegraphics[width=\linewidth]{lecture_02_figure_02.png}
\caption{동일한 점들에 대한 좌표 표기 두 가지와 스칼라장.}
\lectureframe{00:11:25}
\end{figure}

\(x^m\)와 \(y^m\)를 두 개의 매끄러운 좌표계라고 하자. 각각은 점을 순서가 있는 수의 목록으로 보내는 사상이다. 서로 겹치는 좌표 패치에서는 두 목록이 서로 역함수인 다음 함수들로 연결된다.
\begin{equation}
 y^m=y^m(x^1,\ldots,x^d),
 \qquad
 x^m=x^m(y^1,\ldots,y^d).
\label{eq:coordinate-maps}
\end{equation}
매끄러움이 있어야 미분할 수 있고, 가역성이 있어야 좌표를 바꾸는 것만으로 어떤 점도 사라지거나 다른 점과 동일시되지 않는다.

\begin{figure}[H]
\centering
\includegraphics[width=\linewidth]{lecture_02_figure_03.png}
\caption{동일한 기하학적 영역을 기술하는 두 개의 매끄러운 가역 좌표 사상.}
\lectureframe{00:13:10}
\end{figure}

스칼라장의 변환 법칙은 간단히 다음과 같다.
\begin{equation}
 S'(y)=S(x).
\label{eq:scalar-law}
\end{equation}
양변의 함수는 입력값이 다르므로 다르게 보일 수 있지만, 서로 대응하는 점에서의 값은 일치한다. 이것이 불변 명제의 기본형이다.

\section{비스듬한 좌표와 두 종류의 성분}

변환 법칙으로 일반 텐서를 정의하기에 앞서, 벡터의 위 첨자와 아래 첨자를 서로 장식용 표기처럼 취급할 수 없는 이유를 살펴보는 것이 좋다.

축들이 서로 수직도 아니고 정규화되어 있지도 않은 좌표 격자를 그려 보자. 한 점에서 \(x^m\)를 좌표 단위만큼 증가시키면 기저 벡터 \(\mathbf e_m\)가 생긴다. 좌표로 표현하면 다음과 같다.
\begin{equation}
 \mathbf e_m=\frac{\partial \mathbf r}{\partial x^m}.
\label{eq:coordinate-basis}
\end{equation}
벡터 \(\mathbf e_m\)는 길이가 1일 필요도, 서로 직교할 필요도 없다.

\begin{figure}[H]
\centering
\includegraphics[width=\linewidth]{lecture_02_figure_04.png}
\caption{비스듬한 좌표 격자와 그 좌표 기저 벡터.}
\lectureframe{00:17:40}
\end{figure}

모든 벡터 \(\mathbf V\)는 이 기저로 전개할 수 있다.
\begin{equation}
 \mathbf V=V^m\mathbf e_m.
\label{eq:contravariant-expansion}
\end{equation}
수 \(V^m\)는 \emph{반변 성분}이다. 이 성분은 기저 벡터들을 조합해 기하학적 벡터를 만드는 데 필요한 계수다.

\begin{figure}[H]
\centering
\includegraphics[width=\linewidth]{lecture_02_figure_05.png}
\caption{비직교 좌표 기저에서 벡터를 전개한 모습.}
\lectureframe{00:19:20}
\end{figure}

또 다른 자연스러운 수의 모임은 다음과 같다.
\begin{equation}
 V_m=\mathbf V\mathbin{\boldsymbol{\cdot}}\mathbf e_m.
\label{eq:covariant-projection}
\end{equation}
이는 \emph{공변 성분}이다. 비스듬한 기저에서는 단순히 첨자를 옮긴 동일한 계수가 아니다. \eqref{eq:contravariant-expansion}를 대입하면 다음을 얻는다.
\begin{equation}
 V_m
 =V^n\bigl(\mathbf e_n\mathbin{\boldsymbol{\cdot}}\mathbf e_m\bigr).
\label{eq:covariant-from-basis}
\end{equation}

\begin{figure}[H]
\centering
\includegraphics[width=\linewidth]{lecture_02_figure_06.png}
\caption{좌표 기저와의 스칼라곱으로 구한 공변 성분.}
\lectureframe{00:21:10}
\end{figure}

기저 벡터들의 스칼라곱은 그람 행렬을 이룬다.
\begin{equation}
 g_{mn}=\mathbf e_m\mathbin{\boldsymbol{\cdot}}\mathbf e_n.
\label{eq:metric-gram}
\end{equation}
따라서
\begin{equation}
 V_m=g_{mn}V^n.
\label{eq:lower-vector}
\end{equation}
여기서 계량이 처음 등장한다. 계량은 좌표 기저의 길이와 벡터들 사이의 각도를 모두 기록하므로 전개 계수를 사영값으로 바꾸어 준다.

\begin{figure}[H]
\centering
\includegraphics[width=\linewidth]{lecture_02_figure_07.png}
\caption{기저 벡터들의 그람 행렬로 나타낸 계량.}
\lectureframe{00:23:40}
\end{figure}

같은 행렬로 길이도 계산할 수 있다.
\begin{align}
 \mathbf V\mathbin{\boldsymbol{\cdot}}\mathbf V
 &=(V^m\mathbf e_m)\mathbin{\boldsymbol{\cdot}}
   (V^n\mathbf e_n) \notag\\
 &=g_{mn}V^mV^n.
\label{eq:vector-length}
\end{align}

\begin{figure}[H]
\centering
\includegraphics[width=\linewidth]{lecture_02_figure_08.png}
\caption{계량과 반변 성분으로 쓴 벡터의 노름.}
\lectureframe{00:25:00}
\end{figure}

\begin{classroomqa}
\audiencequestion{형식적인 변환 법칙을 정의하기 전에 왜 이런 기하학적 구성을 도입하나요?}
\lecturerresponse{첨자의 의미를 구체적으로 보여 주기 때문이다. 위 첨자 성분은 전개 계수이고, 아래 첨자 성분은 벡터에 대한 선형 측정값이며, 계량은 둘을 연결한다. 그러면 형식적 정의를 통해 이 역할들이 임의의 매끄러운 좌표변환에서도 유지됨을 알 수 있다. \lecturetimestamp{00:26:07}}
\end{classroomqa}

\begin{classroomqa}
\audiencequestion{공변 성분과 반변 성분의 수치가 같아지는 경우는 언제인가요?}
\lecturerresponse{정규직교 유클리드 기저에서는 \(g_{mn}=\delta_{mn}\)이므로, \(V_m=V^m\)가 성분별로 성립합니다. 직교성만으로는 충분하지 않습니다. 기저 벡터의 길이가 단위 길이가 아니면 대각 계량이 여전히 스케일 인자를 제공하기 때문입니다.
\lecturetimestamp{00:26:19}}
\end{classroomqa}

극좌표를 보면 주의할 점이 분명해진다. 반지름 방향과 각도 방향의 좌표선은 서로 직교하지만, 고정된 \(d\theta\)의 변화는 물리적 거리 \(r\,d\theta\)에 해당한다. 평면에서는
\begin{equation}
 ds^2=dr^2+r^2d\theta^2,
\label{eq:polar-metric}
\end{equation}
따라서 이 기저는 직교하지만 데카르트 기저는 아니다.

\begin{figure}[H]
\centering
\includegraphics[width=\linewidth]{lecture_02_figure_09.png}
\caption{극좌표: 반지름에 따라 각도 방향 척도 인자가 커지는 직교 좌표 방향.}
\lectureframe{00:28:10}
\end{figure}

\begin{classroomqa}
\audiencequestion{공변 벡터는 실제로 다른 벡터인가요, 아니면 같은 벡터를 나열하는 또 다른 방법인가요?}
\lecturerresponse{계량을 주기 전에는 벡터와 쌍대벡터는 서로 다른 종류의 선형 대상이다. 벡터는 접공간의 원소이고, 쌍대벡터는 그 공간 위의 선형 함수다. 비퇴화 계량은 둘 사이에 일대일 대응을 정해 주며, 이는 \(V_m=g_{mn}V^n\)로 나타내지만 이 대응을 동일시해서는 안 된다. \lecturetimestamp{00:29:10}}
\end{classroomqa}

\section{변환 법칙으로 정의하는 텐서}

가장 안전한 형식적 정의는 조작적이다. 배열의 성분이 알맞은 야코비 인자들에 따라 변환하면 그 배열은 텐서다. \eqref{eq:coordinate-maps}를 미분하면
\begin{equation}
 dy^m=\frac{\partial y^m}{\partial x^n}\,dx^n.
\label{eq:differential-law}
\end{equation}
모든 반변 벡터는 같은 법칙을 따른다.
\begin{equation}
 V'^m(y)
 =\frac{\partial y^m}{\partial x^n}V^n(x).
\label{eq:contravariant-law}
\end{equation}

\begin{figure}[H]
\centering
\includegraphics[width=\linewidth]{lecture_02_figure_10.png}
\caption{좌표 미분량의 변환을 본보기로 삼은 반변 변환 법칙.}
\lectureframe{00:32:50}
\end{figure}

쌍대벡터는 역야코비 행렬에 따라 변환한다.
\begin{equation}
 W'_m(y)
 =\frac{\partial x^n}{\partial y^m}W_n(x).
\label{eq:covariant-law}
\end{equation}
스칼라의 기울기가 대표적인 예다. \eqref{eq:scalar-law}에 연쇄법칙을 적용하면
\begin{equation}
 \frac{\partial S}{\partial y^m}
 =\frac{\partial x^n}{\partial y^m}
  \frac{\partial S}{\partial x^n}.
\label{eq:gradient-law}
\end{equation}

\begin{figure}[H]
\centering
\includegraphics[width=\linewidth]{lecture_02_figure_11.png}
\caption{연쇄법칙에 따라 쌍대벡터로 변환하는 스칼라 기울기.}
\lectureframe{00:34:20}
\end{figure}

\begin{classroomqa}
\audiencequestion{위 첨자 성분과 아래 첨자 성분은 반드시 물리적 단위가 같나요?}
\lecturerresponse{아니다. 좌표 성분은 좌표와 그 기저 벡터의 단위 및 척도 관례를 물려받는다. 한 좌표가 길이이고 다른 좌표가 각도라면 두 성분은 차원조차 같지 않을 수 있다. 물리적 차원이 좌표계와 무관한 대상은 전체 축약이다.
\lecturetimestamp{00:35:41}}
\end{classroomqa}

\begin{classroomqa}
\audiencequestion{왜 보통 미분이 아니라 편미분을 사용하나요?}
\lecturerresponse{장은 여러 좌표에 의존한다. 야코비 행렬은 한 좌표 함수의 인자 하나만 변하고 나머지는 고정되어 있을 때 그 함수가 어떻게 변하는지를 나타내므로, 그 성분은 편미분이다.
\lecturetimestamp{00:36:45}}
\end{classroomqa}

이제 규칙은 기계적으로 확장된다. 위 첨자가 있는 각 프라임 첨자는 \(\partial y/\partial x\) 인자를 하나 가져오고, 아래 첨자가 있는 각 프라임 첨자는 \(\partial x/\partial y\) 인자를 하나 가져온다. 혼합 텐서는 다음과 같이 변환한다.
\begin{equation}
 T'^m{}_n
 =
 \frac{\partial y^m}{\partial x^p}
 \frac{\partial x^q}{\partial y^n}
 T^p{}_q,
\label{eq:mixed-law}
\end{equation}
랭크 2 공변 텐서는 다음과 같이 변환한다.
\begin{equation}
 W'_{mn}
 =
 \frac{\partial x^p}{\partial y^m}
 \frac{\partial x^q}{\partial y^n}
 W_{pq}.
\label{eq:rank-two-covariant-law}
\end{equation}

\begin{figure}[H]
\centering
\includegraphics[width=\linewidth]{lecture_02_figure_12.png}
\caption{위 첨자 하나와 아래 첨자 하나를 가진 텐서의 변환 법칙.}
\lectureframe{00:40:45}
\end{figure}

\begin{figure}[H]
\centering
\includegraphics[width=\linewidth]{lecture_02_figure_13.png}
\caption{아래 첨자마다 역야코비 인자가 하나씩 들어가는 랭크 2 공변 텐서의 변환 법칙.}
\lectureframe{00:42:15}
\end{figure}

별도 언급이 없으면 반복 첨자에 대해 합산한다. 이 관례는 합 기호를 생략하면서도 이를 복원하는 데 필요한 정보를 모두 보존한다. 더미 첨자는 한 항 안에 정확히 두 번, 한 번은 위 첨자로 한 번은 아래 첨자로 나타나야 한다. 자유 첨자는 방정식의 모든 항에서 한 번씩 나타난다.

\begin{classroomqa}
\audiencequestion{이 합산 표기법은 누가 도입했나요?}
\lecturerresponse{관례상 아인슈타인에게 공로를 돌리므로 아인슈타인 합산 관례라고 부른다. 이 표기법의 가치는 역사적 장식에 있지 않다. 첨자 배열만 보아도 허용되는 텐서 연산을 곧바로 알 수 있다. \lecturetimestamp{00:43:21}}
\end{classroomqa}

\section{텐서 방정식과 좌표계에 무관한 진리}

성분은 좌표계에 따라 달라지지만, 텐서 방정식은 모든 좌표계에서 성립한다. 같은 종류의 두 텐서가 \(W=T\)를 만족한다면 그 차는 영 텐서다.
\begin{equation}
 W^{m\cdots}{}_{\cdots n}
 -T^{m\cdots}{}_{\cdots n}=0.
\label{eq:tensor-equality}
\end{equation}
변환 법칙은 선형이므로 한 좌표계에서 성분이 모두 0이면 모든 좌표계에서도 성분이 모두 0이다. 방정식 양변에서 자유 첨자가 일치해야 한다는 점은 필수적이다. 텐서 종류가 다른 대상 사이의 방정식에는 좌표계에 무관한 의미가 없다.

\begin{figure}[H]
\centering
\includegraphics[width=\linewidth]{lecture_02_figure_14.png}
\caption{첨자 구조가 일치해야 성분의 등식이 모든 좌표계에서 성립하는 텐서 방정식이 된다.}
\lectureframe{00:46:35}
\end{figure}

\begin{classroomqa}
\audiencequestion{벡터 성분이 변하는데 정확히 무엇이 불변인가요?}
\lecturerresponse{기하학적 텐서와 텐서 방정식의 참·거짓은 불변이지만 좌표 성분 목록은 불변이 아니다. 스칼라는 첨자가 하나도 없는 특별한 경우이므로, 한 점에서의 수치는 모든 좌표계에서 같다. \lecturetimestamp{00:47:19}}
\end{classroomqa}

\begin{classroomqa}
\audiencequestion{벡터의 크기가 불변량에 해당하나요? 크기가 0이면 벡터도 0인가요?}
\lecturerresponse{계량이 주어지면 축약 \(g_{mn}V^mV^n\)은 스칼라이므로 불변인 제곱 노름이다. 양의 정부호 기하학에서는 영벡터일 때만 이 값이 0이다. 로런츠 시공간에서는 영이 아닌 널 벡터의 노름도 0일 수 있으므로, 두 명제는 구분해야 한다. \lecturetimestamp{00:48:04}}
\end{classroomqa}

\section{허용되는 텐서 대수}

칠판에는 앞으로 필요한 주요 연산이 정리되어 있다. 스칼라를 곱하고, 같은 종류의 텐서를 더하거나 빼고, 텐서곱을 만들고, 위 첨자와 아래 첨자를 축약하며, 마지막으로 공변 미분을 한다.

\begin{figure}[H]
\centering
\includegraphics[width=\linewidth]{lecture_02_figure_15.png}
\caption{칠판에 정리한 텐서 연산: 덧셈, 곱셈, 축약, 미분.}
\lectureframe{00:50:40}
\end{figure}

\begin{classroomqa}
\audiencequestion{무엇에 대한 미분인가요?}
\lecturerresponse{위치, 따라서 좌표에 대한 미분이다. 성분을 보통 미분한 결과는 일반적으로 텐서처럼 변환하지 않는다. 야코비 행렬 자체가 위치에 따라 변하기 때문이다. 공변 미분은 이를 바로잡지만, 그 구성은 기하학적 기초를 갖춘 뒤로 미룬다. \lecturetimestamp{00:51:49}}
\end{classroomqa}

스칼라 곱셈은 자명하다. 덧셈과 뺄셈은 텐서 종류가 일치해야 한다.
\begin{equation}
 C^m{}_n=A^m{}_n+B^m{}_n,
\qquad
 D^m{}_n=A^m{}_n-B^m{}_n.
\label{eq:tensor-addition}
\end{equation}
벡터와 쌍대벡터를 더하거나 랭크 2 텐서와 스칼라를 더할 수는 없다. 속도와 면적을 더할 수 없는 것과 같다.

\begin{figure}[H]
\centering
\includegraphics[width=\linewidth]{lecture_02_figure_16.png}
\caption{자유 첨자 패턴이 같은 텐서끼리만 성분별 덧셈을 정의할 수 있다.}
\lectureframe{00:53:00}
\end{figure}

텐서곱은 각 인자의 첨자를 모두 결합한다. 예를 들면
\begin{equation}
 T^{mn}=V^mW^n,
\qquad
 S^m{}_n=V^mU_n.
\label{eq:tensor-products}
\end{equation}
4차원에서 \(T^{mn}\)에는 성분이 16개 있다. 어떤 첨자도 합산하지 않았으므로 이는 내적이 아니라 외적이다.

\begin{figure}[H]
\centering
\includegraphics[width=\linewidth]{lecture_02_figure_17.png}
\caption{각 벡터 인자에서 자유 첨자 하나씩을 남기는 텐서곱.}
\lectureframe{00:55:15}
\end{figure}

\begin{classroomqa}
\audiencequestion{여기서 두 벡터의 곱은 내적인가요?}
\lecturerresponse{아니다. 내적은 첨자를 축약해 스칼라를 만든다. 텐서곱은 두 첨자를 모두 자유롭게 남겨 랭크 2 텐서를 만든다. 4차원에서는 성분이 1개가 아니라 16개다. \lecturetimestamp{00:56:36}}
\end{classroomqa}

\begin{classroomqa}
\audiencequestion{성분인 수끼리는 교환할 수 있는데, 두 벡터 인자의 순서를 바꿔도 차이가 없나요?}
\lecturerresponse{보통의 스칼라 성분에 대해서는
\(V^mW^n=W^nV^m\)이다. 그러나 텐서의 각 자리는 순서가 있다. 인자를 서로 바꾼 다음 첫 번째 자리를 \(m\), 두 번째 자리를 \(n\)로 이름 붙이면 전치된 배열이 된다. 따라서 \(T^{mn}\)가 \(T^{nm}\)와 같을 필요는 없다.
\lecturetimestamp{00:58:52}}
\end{classroomqa}

\section{축약과 역야코비 보조정리}

축약의 핵심은 연쇄법칙이다. 좌표 사상들은 서로 역함수이므로
\begin{equation}
 \frac{\partial x^b}{\partial y^m}
 \frac{\partial y^m}{\partial x^a}
 =
 \frac{\partial x^b}{\partial x^a}
 =\delta^b{}_a.
\label{eq:inverse-jacobian}
\end{equation}
두 야코비 행렬은 서로 역행렬이며, 그 곱은 크로네커 델타다.

\begin{figure}[H]
\centering
\includegraphics[width=\linewidth]{lecture_02_figure_18.png}
\caption{텐서 축약의 바탕이 되는 보조정리로서 제시한 역야코비 항등식.}
\lectureframe{01:02:15}
\end{figure}

\begin{classroomqa}
\audiencequestion{역야코비 행렬은 어떤 대상인가요?}
\lecturerresponse{역좌표 사상의 도함수 행렬이다. 그 성분은 \(\partial x^b/\partial y^m\)이며, 다변수 연쇄법칙에 따라 \(\partial y^m/\partial x^a\)와 행렬곱을 하면 항등 행렬이 된다. \lecturetimestamp{01:02:37}}
\end{classroomqa}

간단히 증명해 보면 표기가 분명해진다. 임의의 스칼라 함수 \(f\)에 대해
\begin{equation}
 \frac{\partial f}{\partial y^m}
 \frac{\partial y^m}{\partial x^a}
 =\frac{\partial f}{\partial x^a}.
\end{equation}
\(f=x^b\)를 택하면
\begin{equation}
 \frac{\partial x^b}{\partial y^m}
 \frac{\partial y^m}{\partial x^a}
 =\frac{\partial x^b}{\partial x^a}
 =\delta^b{}_a.
\end{equation}

\begin{figure}[H]
\centering
\includegraphics[width=\linewidth]{lecture_02_figure_19.png}
\caption{연쇄법칙에 따라 역야코비 인자들의 곱이 크로네커 델타로 정리된다.}
\lectureframe{01:04:55}
\end{figure}

\begin{classroomqa}
\audiencequestion{\(\delta^b{}_a\)는 데카르트 좌표에서만 항등 행렬인가요?}
\lecturerresponse{아니다. 혼합 텐서 \(\delta^b{}_a\)는 모든 좌표계에서 항등 사상을 나타낸다. \(\delta_{ab}\)로 쓴 유클리드 계량은 정규직교 데카르트 좌표계에만 특별히 해당하지만, 혼합 첨자를 가진 항등 텐서를 그 계량과 혼동해서는 안 된다.
\lecturetimestamp{01:04:58}}
\end{classroomqa}

이제 혼합 텐서 \(T^m{}_n\)를 생각하자. 이 텐서의 변환된 대각합은 다음과 같다.
\begin{align}
 T'^m{}_m
 &=
 \frac{\partial y^m}{\partial x^a}
 \frac{\partial x^b}{\partial y^m}
 T^a{}_b \notag\\
 &=\delta^b{}_aT^a{}_b
 =T^a{}_a.
\label{eq:trace-invariant}
\end{align}
축약된 양에는 자유 첨자가 없으므로 스칼라다.

\begin{figure}[H]
\centering
\includegraphics[width=\linewidth]{lecture_02_figure_20.png}
\caption{위 첨자와 아래 첨자를 축약하기 전 혼합 텐서의 변환.}
\lectureframe{01:07:45}
\end{figure}

\begin{figure}[H]
\centering
\includegraphics[width=\linewidth]{lecture_02_figure_21.png}
\caption{축약하면 야코비 인자들이 서로 상쇄되어 두 좌표계에서 같은 스칼라 대각합이 남는다.}
\lectureframe{01:10:30}
\end{figure}

\begin{classroomqa}
\audiencequestion{축약한 뒤에도 결과에 첨자가 남나요?}
\lecturerresponse{축약한 한 쌍의 첨자만 사라진다. \(T^m{}_n\)의 유일한 위 첨자와 아래 첨자를 축약하면 스칼라가 남는다. 랭크가 더 높은 텐서에서는 축약하지 않은 첨자들이 모두 자유롭게 남는다.
\lecturetimestamp{01:09:00}}
\end{classroomqa}

예를 들어
\begin{equation}
 A^{mn}{}_{pq}
 \longmapsto
 A^{mn}{}_{pn}
 =B^m{}_p
\label{eq:higher-rank-contraction}
\end{equation}
\(n\)를 축약하여 랭크를 2만큼 낮춘다. 첨자 문자는 더미 표기일 뿐이다. 구조적 규칙은 위 첨자 자리 하나와 아래 첨자 자리 하나를 합산한다는 것이다.

\begin{figure}[H]
\centering
\includegraphics[width=\linewidth]{lecture_02_figure_22.png}
\caption{랭크가 높은 텐서의 축약은 위 첨자 하나와 아래 첨자 하나를 없애면서 나머지 자유 첨자를 보존한다.}
\lectureframe{01:12:35}
\end{figure}

\begin{classroomqa}
\audiencequestion{같은 방식으로 위 첨자 둘이나 아래 첨자 둘을 축약할 수 없는 이유는 무엇인가요?}
\lecturerresponse{위 첨자 둘은 같은 방향의 야코비 인자 둘을 제공한다.
\(\partial y^m/\partial x^a\)와
\(\partial y^m/\partial x^b\)이다. 이 둘은 곱해도 항등 행렬이 되지 않으므로, 제안된 합산은 텐서 변환 법칙을 따르지 않는다. 계량으로 먼저 첨자 하나를 내리면 위 첨자와 아래 첨자의 축약이 가능하다. 그러한 추가 텐서가 없으면 이 연산은 좌표에 불변이지 않다. --
\lecturetimestamp{01:13:03}}
\end{classroomqa}

\begin{figure}[H]
\centering
\includegraphics[width=\linewidth]{lecture_02_figure_23.png}
\caption{같은 종류의 첨자를 축약하려 하면 같은 방향의 야코비 인자 둘이 남으므로 항등 행렬로 정리되지 않는다.}
\lectureframe{01:14:30}
\end{figure}

\begin{classroomqa}
\audiencequestion{축약은 텐서에서 내적에 해당하나요?}
\lecturerresponse{축약은 벡터 자리와 쌍대벡터 자리를 짝짓는 기본적인 불변 쌍이다. 계량이 있으면 두 벡터를 그런 쌍으로 바꿀 수 있으므로, 익숙한 내적은 계량을 적용한 뒤 축약하는 것이다. \lecturetimestamp{01:14:43}}
\end{classroomqa}

\begin{classroomqa}
\audiencequestion{축약하는 두 첨자는 같은 범위에서 변해야 하나요?}
\lecturerresponse{그렇다. 두 첨자는 동일한 \(d\)차원 접공간과 여접공간의 쌍대 기저를 가리킨다. 그렇지 않으면 합산과 역야코비 항등식이 정의되지 않는다.
\lecturetimestamp{01:15:13}}
\end{classroomqa}

\begin{classroomqa}
\audiencequestion{한 기준계에서 영인 텐서는 왜 모든 기준계에서 영인가요?}
\lecturerresponse{변환된 각 성분은 원래 성분들의 선형결합입니다. 입력 성분이 모두 영이면 그 모든 선형결합도 영이 됩니다. 이러한 단순한 선형성 덕분에 텐서 방정식은 좌표계와 무관합니다. \lecturetimestamp{01:15:31}}
\end{classroomqa}

\section{제곱 거리로부터 계량을 정의하기}

이제 다시 기하학으로 돌아가자. 성분이 \(dx^m\)인 무한소 변위의 제곱 길이를 다음과 같이 정의한다.
\begin{equation}
 ds^2=g_{mn}(x)\,dx^m dx^n.
\label{eq:line-element}
\end{equation}
이는 피타고라스 정리를 일반화한 것이다. 계량은 점에 따라 변할 수 있고, 비대각 성분은 좌표 방향들이 서로 직교하지 않음을 나타낸다.

\begin{figure}[H]
\centering
\includegraphics[width=\linewidth]{lecture_02_figure_24.png}
\caption{텐서 연산 목록에서 무한소 선요소의 이차형 정의로 넘어간다.}
\lectureframe{01:18:45}
\end{figure}

\(dx^m dx^n=dx^n dx^m\)이므로 대칭 부분만 기여한다. 따라서 다음과 같이 둘 수 있다.
\begin{equation}
 g_{mn}=g_{nm}.
\label{eq:metric-symmetry}
\end{equation}
\(d\)차원에서 대칭 계량의 독립 성분 수는
\begin{equation}
 \frac{d(d+1)}{2}
\label{eq:metric-component-count}
\end{equation}
이다. 4차원에서는 10개, 3차원에서는 6개, 2차원에서는 3개다.

\begin{figure}[H]
\centering
\includegraphics[width=\linewidth]{lecture_02_figure_25.png}
\caption{선요소와 계량의 독립 성분 수를 세는 데 사용하는 대칭 성분 배열.}
\lectureframe{01:21:35}
\end{figure}

\begin{classroomqa}
\audiencequestion{4차원에서 독립적인 계량 성분이 16개가 아니라 10개인 이유는 무엇인가요?}
\lecturerresponse{대각 성분이 4개이고, 비대각 성분 쌍이 6개다. 다시 말해
\(4(4+1)/2=10\)이다. 반대칭 부분은 대칭곱 \(dx^m dx^n\)에 곱해지므로 상쇄된다. \lecturetimestamp{01:19:50}}
\end{classroomqa}

\begin{classroomqa}
\audiencequestion{3차원에는 독립 성분이 몇 개 있나요?}
\lecturerresponse{6개다. 대각 성분 3개와 비대각 성분 쌍 3개이며, 또는 \(3(3+1)/2\)이다. \lecturetimestamp{01:20:47}}
\end{classroomqa}

\begin{classroomqa}
\audiencequestion{그럼 2차원에서는요?}
\lecturerresponse{3개다. \(g_{11}\), \(g_{22}\), 그리고 대칭 성분 한 쌍 \(g_{12}=g_{21}\)이다. \lecturetimestamp{01:22:28}}
\end{classroomqa}

\subsection{계량이 텐서인 이유}

제곱 거리는 스칼라다. 따라서 같은 변위를 두 좌표계로 쓰면
\begin{equation}
 g_{mn}(x)\,dx^m dx^n
 =
 g'_{pq}(y)\,dy^pdy^q.
\label{eq:length-invariance}
\end{equation}
다음 식을
\begin{equation}
 dx^m=\frac{\partial x^m}{\partial y^p}\,dy^p
\end{equation}
두 번 적용하면
\begin{align}
 ds^2
 &=
 g_{mn}(x)
 \frac{\partial x^m}{\partial y^p}
 \frac{\partial x^n}{\partial y^q}
 dy^pdy^q.
\label{eq:metric-transform-derivation}
\end{align}
프라임 좌표 표현과 비교하면 다음 관계를 얻는다.
\begin{equation}
 g'_{pq}(y)
 =
 \frac{\partial x^m}{\partial y^p}
 \frac{\partial x^n}{\partial y^q}
 g_{mn}(x).
\label{eq:metric-transformation}
\end{equation}
이는 랭크 2 공변 텐서의 변환 법칙과 정확히 일치한다.

\begin{figure}[H]
\centering
\includegraphics[width=\linewidth]{lecture_02_figure_26.png}
\caption{스칼라 선요소의 불변성으로부터 계량 변환 법칙을 유도한다.}
\lectureframe{01:26:25}
\end{figure}

이 결론은 계량에 아래 첨자가 두 개 있는 이유도 설명한다. 아래에 있는 각 계량 첨자는 변위 성분의 위 첨자 하나와 짝을 이룬다. 좌표를 바꿀 때 \eqref{eq:metric-transformation}에 있는 역야코비 인자 두 개는 미분량의 순야코비 인자 두 개와 상쇄되어 \(ds^2\)를 변하지 않게 한다.

\section{계량 행렬과 그 역행렬}

한 점에서 \(g_{mn}\)는 대칭 행렬이다.
\begin{equation}
 [g_{mn}]
 =
 \begin{pmatrix}
 g_{11} & g_{12} & \cdots \\
 g_{12} & g_{22} & \cdots \\
 \vdots & \vdots & \ddots
 \end{pmatrix}.
\label{eq:metric-matrix}
\end{equation}

\begin{figure}[H]
\centering
\includegraphics[width=\linewidth]{lecture_02_figure_27.png}
\caption{텐서 변환 법칙을 확립한 뒤 대칭 행렬로 나타낸 계량.}
\lectureframe{01:28:50}
\end{figure}

계량은 비퇴화여야 한다. 즉, 고윳값이 0인 것이 없으므로 역행렬을 가진다. 그 역행렬을 \(g^{mn}\)로 표기하자. 첨자 표기법으로는
\begin{equation}
 g_{mn}g^{np}=\delta_m{}^p.
\label{eq:inverse-metric}
\end{equation}
반복되는 \(n\)에 대해 합산하고, \(m\)과 \(p\)는 곱의 행과 열을 표시한다.

\begin{figure}[H]
\centering
\includegraphics[width=\linewidth]{lecture_02_figure_28.png}
\caption{역계량은 다음 관계를 통해 도입한다.
\(g_{mn}g^{np}=\delta_m{}^p\)}
\lectureframe{01:31:40}
\end{figure}

\begin{classroomqa}
\audiencequestion{한 점에서 벡터와 텐서는 실제로 어디에 놓이나요?}
\lecturerresponse{한 점의 벡터들은 그 점의 접공간 \(T_P M\)을 이루고, 쌍대벡터들은 쌍대공간 \(T_P^*M\)을 이룬다. 고차 텐서는 이 공간들의 텐서곱에 속한다. 텐서장은 각 점에 그러한 대상을 매끄럽게 하나씩 대응시킨다. 계량은 접공간 위의 쌍선형 형식이고, 역계량은 여접공간 위의 대응하는 쌍선형 형식이다. \lecturetimestamp{01:34:15}}
\end{classroomqa}

\begin{figure}[H]
\centering
\includegraphics[width=\linewidth]{lecture_02_figure_29.png}
\caption{곡률 다양체의 한 점에 붙은 접공간에 그린 벡터.}
\lectureframe{01:35:45}
\end{figure}

\begin{classroomqa}
\audiencequestion{\(g_{mn}g^{np}=\delta_m{}^p\)는 \(m=p\)인 경우만 가리키나요?}
\lecturerresponse{아니다. 완전한 행렬 방정식이다. 모든 순서쌍 \((m,p)\)에 대해, 합산된 곱은 \(m=p\)이면 1이고 \(m\ne p\)이면 0이다. 대각 방정식과 비대각 방정식이 모두 동시에 들어 있다. \lecturetimestamp{01:36:50}}
\end{classroomqa}

\begin{figure}[H]
\centering
\includegraphics[width=\linewidth]{lecture_02_figure_30.png}
\caption{항등 행렬 전체를 나타내는 혼합 첨자 관계로 해석한 역관계.}
\lectureframe{01:38:50}
\end{figure}

2차원에서는 다음을 뜻한다.
\begin{equation}
 \begin{pmatrix}
 g_{11} & g_{12}\\
 g_{12} & g_{22}
 \end{pmatrix}
 \begin{pmatrix}
 g^{11} & g^{12}\\
 g^{12} & g^{22}
 \end{pmatrix}
 =
 \begin{pmatrix}
 1&0\\
 0&1
 \end{pmatrix}.
\label{eq:two-dimensional-inverse}
\end{equation}
예를 들면
\begin{align}
 g_{11}g^{11}+g_{12}g^{12}&=1,\\
 g_{11}g^{12}+g_{12}g^{22}&=0.
\end{align}
다른 행에서는 나머지 두 방정식을 얻는다.

\begin{figure}[H]
\centering
\includegraphics[width=\linewidth]{lecture_02_figure_31.png}
\caption{2차원 행렬곱으로 명시적으로 전개한 역계량 항등식.}
\lectureframe{01:40:00}
\end{figure}

\begin{classroomqa}
\audiencequestion{계량은 양의 정부호여야 하나요?}
\lecturerresponse{보통의 리만 기하학에서는 양의 정부호성이 모든 영이 아닌 벡터의 제곱 길이를 양수로 만든다. 시공간에는 로런츠 계량을 사용하며, 부호 하나가 나머지 세 부호와 반대다. 여기서 두 경우에 공통으로 필요한 것은 양의 정부호성이 아니라 비퇴화성이다.
\lecturetimestamp{01:42:18}}
\end{classroomqa}

\begin{classroomqa}
\audiencequestion{고윳값이 0인 경우는 왜 제외하나요?}
\lecturerresponse{고윳값이 0이면 계량 행렬이 특이행렬이 되어 \(g^{mn}\)가 존재하지 않으며, 계량은 벡터와 공변벡터 사이에 일대일 대응을 정할 수 없습니다. 통상적인 미분기하학적 의미의 계량은 비퇴화여야 합니다. \lecturetimestamp{01:43:56}}
\end{classroomqa}

\begin{classroomqa}
\audiencequestion{4차원 계량의 독립적인 성분 10개는 10차원 끈 이론과 관련이 있나요?}
\lecturerresponse{아닙니다. 앞의 10은 단지 대칭인 4×4 행렬의 성분 수인 \(4(4+1)/2\)입니다. 초끈 이론에서 시공간 차원이 10으로 나타나는 데에는 다른 기원이 있습니다. \lecturetimestamp{01:44:32}}
\end{classroomqa}

\begin{classroomqa}
\audiencequestion{로런츠 시공간에 길이가 0인 영벡터가 있다면, 계량은 역행렬을 갖지 못하게 되나요?}
\lecturerresponse{아닙니다. 영벡터는 \(g_{mn}V^mV^n=0\)를 만족하는데, 이는 부정부호 이차 형식 안에서 항들이 상쇄된 결과입니다. 이는 \(g_{mn}V^n=0\)을 만족하는 고윳값 0의 벡터와는 다릅니다. 예를 들어, 민코프스키 계량의 고윳값은 \(\pm1\)이며, 빛꼴 벡터를 허용하면서도 역행렬을 완벽하게 갖습니다. \lecturetimestamp{01:45:02}}
\end{classroomqa}

\section{다음에 다룰 내용}

이제 기하학적 구상을 명료하게 서술할 수 있습니다. 계량은 무한소 거리를 나타내고, 공변적으로 변환되며, 역을 갖습니다. 좌표가 복잡하기만 해도 계량의 성분은 복잡해 보일 수 있습니다. 기하학 자체가 휘어 있는지 판단하려면, 좌표에 무관한 방식으로 이웃한 점들에서의 계량을 비교해야 합니다.

텐서를 미분하면 변환 법칙에 들어 있는 위치 의존 야코비안도 함께 미분되므로, 보통의 미분만으로는 충분하지 않습니다. 따라서 다음 단계는 공변미분입니다. 계량의 미분으로부터 접속을 구성하고, 그 접속의 미분과 곱으로부터 곡률을 구성합니다. 그래야만 말린 페이지에 관한 질문에 기하학 자체에 내재한 최종 답을 얻을 수 있습니다.
